\chapter{Ausgewählte klassische arithmetische Hauptsätze}
\label{b10guide:ch:classical}

Die folgenden Sätze ergänzen die grundlegenden Rechenregeln um
klassische Aussagen über natürliche Zahlen. Alle Variablen und
alle gesuchten Lösungen liegen in $\mathbb N$. Restvorgaben werden
mit dem bereits definierten natürlichen Restwert formuliert.

\section{Potenzen, Fakultäten und Restvorgaben}

\begin{theorem}[Legendre-Formel für Fakultäten]
\label{b10guide:thm:legendre}
Für jede Primzahl $p$ und $n\in\mathbb N$ gilt
\[
 v_p(n!)=\sum_{j<n}
       \bigl(n\mathbin{\mathrm{div}}p^{j+1}\bigr).
\]
Der Fall $n=0$ hat auf beiden Seiten den Wert $0$.
Insbesondere gilt für Primzahlen $p$
\[
 p\mid n!\quad\Longleftrightarrow\quad p\le n.
\]
Vergleichsquelle: \cite{natlegendre}.
\end{theorem}

\begin{theorem}[Fermats kleiner Satz]
\label{b10guide:thm:fermat-little}
Für jede Primzahl $p$ und jede natürliche Zahl $a$ gilt
\[
 a^p\mathbin{\mathrm{mod}}p=a\mathbin{\mathrm{mod}}p.
\]
Falls $p\nmid a$, gilt außerdem
\[
 a^{p-1}\mathbin{\mathrm{mod}}p=1.
\]
Vergleichsquelle: \cite{natfermat}.
\end{theorem}

\begin{theorem}[Wilson-Kriterium für Primzahlen]
\label{b10guide:thm:wilson}
Für jede natürliche Zahl $n\ge2$ gilt
\[
 n\text{ ist prim}\quad\Longleftrightarrow\quad
 n\mid\bigl((n-1)!+1\bigr).
\]
Gleichwertig ist die Bedingung
$(n-1)!\mathbin{\mathrm{mod}}n=n-1$ \cite{natwilson}.
\end{theorem}

\begin{theorem}[Chinesischer Restsatz für zwei natürliche Divisoren]
\label{b10guide:thm:chinese-remainder}
Seien $m,n\in\mathbb N_+$ mit $\operatorname{ggT}(m,n)=1$
und seien $r,s\in\mathbb N$ mit $r<m$ und $s<n$.
Dann gibt es genau ein $x\in\mathbb N$ mit
\[
 x<mn,\qquad x\mathbin{\mathrm{mod}}m=r,\qquad
 x\mathbin{\mathrm{mod}}n=s.
\]
Alle natürlichen Zahlen mit diesen beiden Restwerten sind genau
\[
 y=x+kmn\qquad(k\in\mathbb N).
\]
Dabei ist $k$ für jede Lösung $y$ eindeutig \cite{natchinese}.
\end{theorem}

\section{Primzahlen und zusammengesetzte Zahlen zwischen Schranken}

\begin{theorem}[Bertrandscher Satz]
\label{b10guide:thm:bertrand}
Für jede natürliche Zahl $n\ge2$ gibt es eine Primzahl $p$ mit
\[
 n<p<2n.
\]
Allgemeiner gibt es für jedes $n\ge1$ eine Primzahl $p$ mit
$n<p\le2n$ \cite{natbertrand}.
\end{theorem}

\begin{theorem}[Beliebig lange Blöcke zusammengesetzter Zahlen]
\label{b10guide:thm:composite-blocks}
Zu jeder Länge $k\in\mathbb N_+$ und jeder Schranke $B\in\mathbb N$
gibt es eine Zahl $a>B$, für die alle Zahlen
\[
 a,\ a+1,\ \ldots,\ a+(k-1)
\]
zusammengesetzt sind.
\end{theorem}

\section{Darstellung natürlicher Zahlen durch Quadrate}

\begin{theorem}[Zwei-Quadrate-Kriterium]
\label{b10guide:thm:two-squares}
Eine positive natürliche Zahl $n$ besitzt genau dann eine Darstellung
$n=a^2+b^2$ mit $a,b\in\mathbb N$, wenn für jede Primzahl $p$ mit
$p\mathbin{\mathrm{mod}}4=3$ der Primexponent $v_p(n)$ gerade ist.
Insbesondere ist eine Primzahl $p$ genau dann Summe zweier
natürlicher Quadrate, wenn $p=2$ oder
$p\mathbin{\mathrm{mod}}4=1$ gilt.
Für die Null gilt gesondert $0=0^2+0^2$ \cite{natsquares}.
\end{theorem}

\begin{theorem}[Vier-Quadrate-Satz von Lagrange]
\label{b10guide:thm:four-squares}
Für jede natürliche Zahl $n$ gibt es $a,b,c,d\in\mathbb N$ mit
\[
 n=a^2+b^2+c^2+d^2.
\]
Nullen und wiederholte Werte sind zugelassen; eine Eindeutigkeit
der Darstellung wird nicht vorausgesetzt \cite{natsquares}.
\end{theorem}

\section{Weitere Restgleichungen}

\begin{theorem}[Allgemeiner chinesischer Restsatz für zwei Divisoren]
\label{b10guide:thm:general-chinese-remainder}
Seien $m,n\in\mathbb N_+$ und $r,s\in\mathbb N$ mit $r<m$ und $s<n$.
Setze $d=\operatorname{ggT}(m,n)$ und $L=\operatorname{kgV}(m,n)$.
Es gibt genau dann eine natürliche Zahl $x$ mit
\[
 x\mathbin{\mathrm{mod}}m=r,\qquad
 x\mathbin{\mathrm{mod}}n=s,
\]
wenn
\[
 r\mathbin{\mathrm{mod}}d=s\mathbin{\mathrm{mod}}d.
\]
Ist diese Bedingung erfüllt, gibt es genau eine solche Zahl $x_0<L$.
Alle natürlichen Lösungen sind dann genau
\[
 x=x_0+kL\qquad(k\in\mathbb N),
\]
wobei $k$ für jede Lösung eindeutig ist \cite{natchinese}.
\end{theorem}

\begin{theorem}[Lösbarkeit und sämtliche Lösungen einer linearen Restgleichung]
\label{b10guide:thm:linear-remainder-equation}
Seien $a,b\in\mathbb N$ und $m\in\mathbb N_+$.
Setze $d=\operatorname{ggT}(a,m)$ und
$h=m\mathbin{\mathrm{div}}d$; dabei sind $d$ und $h$ positiv.
Die Restgleichung
\[
 (ax)\mathbin{\mathrm{mod}}m=b\mathbin{\mathrm{mod}}m
\]
hat genau dann eine Lösung $x\in\mathbb N$, wenn $d\mid b$ gilt.
In diesem Fall gibt es genau eine Lösung $x_0<h$, und sämtliche
natürlichen Lösungen sind
\[
 x=x_0+kh\qquad(k\in\mathbb N).
\]
Der Parameter $k$ ist für jede Lösung eindeutig.
Die Aussage gilt auch für $a=0$ und für $m=1$ \cite{natlinearrest}.
\end{theorem}

\section{Weitere natürliche Zahlendarstellungen}

\begin{theorem}[Frobenius-Satz für zwei teilerfremde natürliche Erzeuger]
\label{b10guide:thm:frobenius-two-generators}
Seien $a,b\in\mathbb N$ mit $a\ge2$, $b\ge2$ und
$\operatorname{ggT}(a,b)=1$.
Wegen $a+b\le ab$ ist die natürliche Differenz
\[
 g=ab-(a+b)
\]
definiert. Es gibt keine $x,y\in\mathbb N$ mit $g=ax+by$.
Für jedes $N\in\mathbb N$ mit $N>g$ gibt es hingegen
$x,y\in\mathbb N$ mit
\[
 N=ax+by.
\]
Damit ist $g$ die größte natürliche Zahl, die keine solche Darstellung
besitzt \cite{natfrobenius}.
\end{theorem}

\begin{theorem}[Parametrisierung pythagoreischer Tripel]
\label{b10guide:thm:pythagorean-triples}
Seien $a,b,c\in\mathbb N_+$.
Die Gleichung $a^2+b^2=c^2$ gilt genau dann, wenn es
$d,u,v\in\mathbb N_+$ mit
\[
 u>v,\qquad \operatorname{ggT}(u,v)=1
\]
und verschiedener Parität von $u$ und $v$ gibt, sodass nach
gegebenenfalls erforderlicher Vertauschung von $a$ und $b$ gilt:
\[
 a=d(u^2-v^2),\qquad b=2duv,\qquad c=d(u^2+v^2).
\]
Dabei ist $d=\operatorname{ggT}(a,b)$.
Die natürlichen Lösungen mit einer Null unter $a,b,c$ sind genau
$(0,t,t)$ und $(t,0,t)$ mit $t\in\mathbb N$ \cite{natpythagoras}.
\end{theorem}

\begin{theorem}[Drei-Quadrate-Kriterium von Legendre]
\label{b10guide:thm:three-squares}
Eine natürliche Zahl $n$ besitzt genau dann eine Darstellung
\[
 n=x^2+y^2+z^2\qquad(x,y,z\in\mathbb N),
\]
wenn es keine $a,b\in\mathbb N$ mit
\[
 n=4^a(8b+7)
\]
gibt. Nullen unter den Summanden sind zugelassen; insbesondere
ist $0=0^2+0^2+0^2$ eingeschlossen \cite{natthreesquares}.
\end{theorem}


\section{Primzahlen mit vorgegebenem Rest}

\begin{theorem}[Dirichlets Satz über Primzahlen mit vorgegebenem Rest]
\label{b10guide:thm:dirichlet-prime-remainders}
Seien $q\in\mathbb N_+$ und $a\in\mathbb N$ mit
$\operatorname{ggT}(a,q)=1$. Zu jeder natürlichen Schranke $B$
gibt es eine Primzahl $p$ mit
\[
 p>B,\qquad p\mathbin{\mathrm{mod}}q
                 =a\mathbin{\mathrm{mod}}q.
\]
Die Primzahlen mit diesem Rest sind somit unbeschränkt.
Auch der Fall $q=1$ ist eingeschlossen \cite{natdirichlet}.
\end{theorem}

\section{Höhere Potenzgleichungen}

\begin{theorem}[Fermats letzter Satz]
\label{b10guide:thm:fermat-last}
Für jede natürliche Zahl $r\ge3$ gibt es keine positiven natürlichen
Zahlen $x,y,z$ mit
\[
 x^r+y^r=z^r.
\]
Gleichwertig gilt für $x,y,z\in\mathbb N$ und $r\ge3$:
Aus $x^r+y^r=z^r$ folgt $xy=0$ \cite{natfermatlast}.
\end{theorem}
